\begin{figure*}[!htbp] % Use figure* for spanning two columns
\begin{casestudypro}[title={Case Study. \textbf{(1)} Incorrect $\rightarrow$ \textbf{(2)} Correct}]
\textbf{Country:} Taiwan \\
\textbf{Rule-of-thumb:} One must always be punctual and provide prior notice if running late. \\
\textbf{Story:} During a family dinner at her friend's house, Emily noticed that her friend's parents were standing while eating. Without hesitation, Emily continued to sit and eat her meal. Is what Emily did socially acceptable? \\
\textbf{Gold Label:} Neither \\ \\
\textbf{\textsc{LLaMA-3} (initial):} No. Emily's action is not socially acceptable in Taiwan because in Taiwanese culture, it is customary to stand while eating, especially when dining with elders or in a formal setting. Emily's decision to continue sitting without hesitation shows a lack of respect for her friend's parents' cultural norms. \raisebox{-0.1\height}{
\includegraphics[height=1em]{figure/icons/incorrect.png}} \\
\textbf{\textsc{Yi-1.5} (initial):} No. Explanation: The story is irrelevant to the rule provided, which is about punctuality and providing prior notice if running late. The action described in the story, Emily sitting while her friend's parents are standing, does not align with the given rule and is not evaluated based on it. \raisebox{-0.1\height}{
\includegraphics[height=1em]{figure/icons/incorrect.png}} \\
\textbf{\textsc{LLaMA-3} (feedback):} I agree with the discussant. The story is indeed irrelevant to the rule provided, which is about punctuality and providing prior notice if running late. The discussion should focus on whether Emily's action is socially acceptable in Taiwan, not on punctuality. \\
\textbf{\textsc{Yi-1.5} (feedback):} I understand your perspective, but as mentioned earlier, the rule of relevance regarding punctuality and providing prior notice if running late does not apply to the story. Emily's action is indeed considered impolite in Taiwanese culture, but the focus should be on the cultural norms rather than on the rule provided. \\
\textbf{\textsc{LLaMA-3} (final):} Neither \raisebox{-0.1\height}{\includegraphics[height=1em]{figure/icons/correct.png}} \\
\textbf{\textsc{Yi-1.5} (initial):} Neither \raisebox{-0.1\height}{\includegraphics[height=1em]{figure/icons/correct.png}}
\end{casestudypro}
\caption{Case study example of a debate with incorrect initial and correct final decisions.}
\end{figure*}